\documentclass[11pt]{article}

\usepackage[a4paper,margin=1in]{geometry}
\usepackage[T1]{fontenc}
\usepackage{lmodern}
\usepackage{microtype}
\usepackage{amsmath,amssymb,amsthm,mathtools}
\usepackage{booktabs,tabularx,array}
\usepackage{enumitem}
\usepackage[numbers,sort&compress]{natbib}
\usepackage[colorlinks=true,linkcolor=blue,citecolor=blue,urlcolor=blue]{hyperref}
\usepackage[nameinlink,noabbrev]{cleveref}

\setlength{\emergencystretch}{3em}

\newtheorem{theorem}{Theorem}[section]
\newtheorem{proposition}[theorem]{Proposition}
\newtheorem{corollary}[theorem]{Corollary}
\newtheorem{lemma}[theorem]{Lemma}
\theoremstyle{definition}
\newtheorem{definition}[theorem]{Definition}
\theoremstyle{remark}
\newtheorem{remark}[theorem]{Remark}

\newcommand{\B}{\mathcal B}
\newcommand{\Cpre}{\mathcal C}
\newcommand{\Ttail}{\mathcal T}
\newcommand{\one}{\mathbf 1}
\newcommand{\Lzero}{\mathrm{L0}}
\newcommand{\Lone}{\mathrm{L1}}
\newcommand{\Ltwo}{\mathrm{L2}}
\newcommand{\NT}{\textsc{not-testable}}
\newcommand{\STOP}{\textsc{stop-declared-route}}
\newcolumntype{P}[1]{>{\raggedright\arraybackslash}p{#1}}
\newcolumntype{Y}{>{\raggedright\arraybackslash}X}

\title{\textbf{A Boundary-Complete Dyadic Prefix--Tail Path Archive}\\
\large Exact Column Conservation and Fixed-Shift Provenance}
\author{Liang Wang}
\date{July 2026}

\begin{document}
\maketitle

\begin{abstract}
TPC-133 supplies an executable native entrance but stops before the
first block decomposition.  We construct a complete growing path
archive through the dyadic and divisor-prefix operations used in
TPC-16 and TPC-17.  An explicit normalized \(C^\infty\) bump gives
\(\sum_j\psi(x/2^j)=1\) exactly.  For every dyadic block \(b\), a
predeclared integer \(D_0(b)\in[0,V]\) sends each opened divisor
record to exactly one of \(d\le D_0(b)\) and
\(D_0(b)<d\le V\).  Thus every native column has outgoing weight one,
the complete matrix satisfies
\[
 \one^TM_{134}=\one^T,
\]
and native provenance, the prescribed shift, and the single physical
normalization intertwine on every path.  The resulting identity
\[
 B_{h_0,\delta}(X)=
 \sum_b\bigl(\mathcal C_b(D_0(b))+
                  \mathcal T_b(D_0(b),V)\bigr)
\]
is exact for every \(X\).  No prefix or tail estimate is inferred.
The partition theorem is \(\Lzero\), its literal attachment is
\(\Lone\), and no positive \(\Ltwo\) arithmetic saving is proved.
\end{abstract}

\section{Input archive and an explicit smooth partition}

Import the support-envelope native records
\[
 \alpha=(\ell,k,d),\qquad
 c_X(\alpha)=
 -\Lambda(\ell)\mu(d)r_Q(\ell k+h_0)W(\ell k/X)
\]
from TPC-133 \citep{WangTPC133}.  They retain \(d\) rather than only
the compressed pair \((\ell,k)\).

Define
\[
 \chi(x)=
 \begin{cases}
 \exp\!\left(-\dfrac1{x-1/2}-\dfrac1{2-x}\right),
   &1/2<x<2,\\
 0,&\text{otherwise},
 \end{cases}
\]
and
\[
 Z(x)=\sum_{j\in\mathbb Z}\chi(x/2^j),
 \qquad
 \psi(x)=\frac{\chi(x)}{Z(x)}.
\]
For every \(x>0\), the sum defining \(Z(x)\) is finite and positive.
Moreover \(Z(x/2^j)=Z(x)\), by reindexing.

\begin{lemma}[Exact normalized dyadic partition]
\label{lem:partition}
The function \(\psi\) is smooth, compactly supported in
\([1/2,2]\), and
\[
 \boxed{\sum_{j\in\mathbb Z}\psi(x/2^j)=1}
 \qquad(x>0).
\]
Only finitely many summands are nonzero for each \(x\).
\end{lemma}

\begin{proof}
Smoothness at the endpoints is the standard flat-bump property.  The
support gives local finiteness.  Since all denominators equal \(Z(x)\),
\[
 \sum_j\psi(x/2^j)
 =\frac{\sum_j\chi(x/2^j)}{Z(x)}=1.
\]
\end{proof}

This explicit formula removes an ambiguity in a machine archive:
the program stores a symbolic \(\psi\)-node and the finite dyadic
orbit.  It does not approximate the column sum in binary floating
point.

\section{The prefix--tail compiler}

Let \(\B_X\) be the finite set of dyadic pairs
\[
 b=(j_L,j_K),\qquad L=2^{j_L},\quad K=2^{j_K},
\]
which meet at least one native atom.  Fix before inspecting any
coefficient a deterministic policy
\[
 D_0:\B_X\longrightarrow\{0,1,\ldots,V\}.
\]
The exact theorem below works for every such policy.  The bundled
certificate uses the maximal integer satisfying a separately declared
published-core screen, with value zero when the screen is empty.

\begin{definition}[Path children]
For \(\alpha=(\ell,k,d)\) and \(b=(j_L,j_K)\), create one child when
\[
 \psi(\ell/L)\psi(k/K)\ne0.
\]
Its edge multiplier is
\[
 m(\alpha,b)=\psi(\ell/L)\psi(k/K),
\]
and its terminal type is
\[
 \operatorname{type}(\alpha,b)=
 \begin{cases}
 \mathrm{PREFIX},&d\le D_0(b),\\
 \mathrm{TAIL},&D_0(b)<d\le V.
 \end{cases}
\]
\end{definition}

\begin{theorem}[Boundary-complete path conservation]
\label{thm:archive}
For every native atom,
\[
 \sum_{b\in\B_X}m(\alpha,b)=1.
\]
The prefix and tail types are disjoint and exhaustive, including the
boundary \(d=D_0(b)\).  Consequently the complete path matrix
\(M_{134}\) satisfies
\[
 \boxed{\one^TM_{134}=\one^T.}
\]
\end{theorem}

\begin{proof}
By \cref{lem:partition},
\[
 \sum_{j_L,j_K}
 \psi(\ell/2^{j_L})\psi(k/2^{j_K})
 =
 \left(\sum_{j_L}\psi(\ell/2^{j_L})\right)
 \left(\sum_{j_K}\psi(k/2^{j_K})\right)=1.
\]
Compact support makes the sum finite.  For each block and each
\(d\le V\), exactly one of \(d\le D_0(b)\) and \(d>D_0(b)\) holds.
Thus no second multiplier is introduced by the type split.
\end{proof}

\begin{theorem}[Metadata intertwining]
\label{thm:metadata}
Suppose every child copies its native tuple, \(h_0\), and the inherited
physical normalization.  Then each of these coordinate projections
intertwines with \(M_{134}\).  In particular no averaged shift or
second normalization is created by the block compiler.
\end{theorem}

\begin{proof}
Every nonzero edge has equal source and target metadata.  The
entrywise criterion of TPC-123 applies \citep{WangTPC123}.
\end{proof}

\section{Exact scalar block identity}

For a block \(b\), define
\begin{align*}
 \Cpre_b(D_0(b))
 &=
 \sum_{\ell,k}
 \sum_{\substack{d\mid k\\d\le D_0(b)}}
 c_X(\ell,k,d)\psi(\ell/L)\psi(k/K),\\
 \Ttail_b(D_0(b),V)
 &=
 \sum_{\ell,k}
 \sum_{\substack{d\mid k\\D_0(b)<d\le V}}
 c_X(\ell,k,d)\psi(\ell/L)\psi(k/K).
\end{align*}

\begin{corollary}[Growing prefix--tail archive identity]
\label{cor:identity}
For every \(X\),
\[
 \boxed{
 B_{h_0,\delta}(X)=
 \sum_{b\in\B_X}
 \bigl(\Cpre_b(D_0(b))+\Ttail_b(D_0(b),V)\bigr).}
\]
\end{corollary}

\begin{proof}
Apply \cref{thm:archive} to the native coefficient vector and then
partition the opened divisors at \(D_0(b)\).  This is the exact
coefficient split used in TPC-17 \citep{WangTPC16,WangTPC17}.
\end{proof}

\begin{remark}[No parameter-transport shortcut]
TPC-16 also proves that two complete Vaughan hard packets can differ
by \(O_A(X(\log X)^{-A})\).  It explicitly warns that this total signed
packet equivalence does not match individual factor pairs.  We do not
insert it as a native-column bijection.  A proof-carrying transport
would have to archive its soft row terms separately.
\end{remark}

\section{Executable certificate and limits}

The standard-library program
\texttt{experiments/tpc134\_branch\_archive.py} reads the committed
TPC-133 JSONL records, computes every finite dyadic orbit, applies one
canonical \(D_0\) policy, and writes a path JSONL plus a summary
certificate.  It rejects a deleted or duplicated overlap child, a
prefix/tail boundary mutation, a wrong source, a wrong shift, and a
second normalization.  The certificate checks the exact finite index
set and records column conservation as an invocation of
\cref{lem:partition}; it does not replace that symbolic identity by a
floating-point column-sum test.  Its \texttt{--check} mode is
read-only.

\begin{center}
\small
\begin{tabularx}{\textwidth}{@{}P{0.43\textwidth}Y@{}}
\toprule
Statement & Status\\
\midrule
Smooth dyadic identity and divisor split
& Exact \(\Lzero\).\\
Complete native-to-block path archive
& Proved at this cut, \(\Lone\).\\
Uniform smallness of prefix blocks
& Not invoked in this paper.\\
Uniform smallness of tail blocks
& Not proved.\\
Physical/determinant/zero-mode total maps
& Not supplied.\\
Positive fixed-\(h_0\) arithmetic power saving
& Not proved; no positive \(\Ltwo\).\\
\bottomrule
\end{tabularx}
\end{center}

\paragraph{Kill criteria.}
A nonzero column defect, a gap or overlap at \(d=D_0\), a lost source,
or a scope mutation is \(\STOP\) for the declared compiler.  A missing
later carrier map is \(\NT\), not a contradiction.

\paragraph{Explicit exclusions.}
This paper proves no determinant-energy lower bound, zero-mode
estimate, physical cover, signed affine cancellation, endpoint below
\(1/400\), Hardy--Littlewood asymptotic, prime-pair lower bound, or
twin-prime theorem.

\begingroup
\small
\setlength{\bibsep}{2pt plus 0.3ex}
\bibliographystyle{plainnat}
\bibliography{references}
\endgroup

\end{document}
